
\documentclass{article}
\usepackage{colortbl}
\usepackage{makecell}
\usepackage{multirow}
\usepackage{supertabular}

\begin{document}

\newcounter{utterance}

\centering \large Interaction Transcript for game `cladder', experiment `full\_v1.5\_default', episode 9908 with Qwen3.6{-}27B.
\vspace{24pt}

{ \footnotesize  \setcounter{utterance}{1}
\setlength{\tabcolsep}{0pt}
\begin{supertabular}{c@{$\;$}|p{.15\linewidth}@{}p{.15\linewidth}p{.15\linewidth}p{.15\linewidth}p{.15\linewidth}p{.15\linewidth}}
    \# & \multicolumn{2}{c}{Player} && \multicolumn{2}{c}{Game Master} \\
    \hline

    \theutterance \stepcounter{utterance}  
    & & & \multicolumn{4}{p{0.6\linewidth}}{
        \cellcolor[rgb]{0.9,0.9,0.9}{
            \makecell[{{p{\linewidth}}}]{
                \texttt{\tiny{[P1$\langle$GM]}}
                \texttt{You are an expert in causal inference. The following question is not a typical commonsense query, but rather a meticulously designed question created by a professor specializing in causal inference, intended to assess the students' mastery of the course content.} \\
\texttt{Imagine a self{-}contained, hypothetical world with only the following conditions, and without any unmentioned factors or causal relationships: Zuph has a direct effect on wibl. Wibl has a direct effect on uvzi. We know that zuph causes not wibl, and we know that wibl causes not uvzi. Would an individual is not uvzi if zuph instead of not zuph?} \\
\texttt{Start your answer with 'Yes' or 'No', followed by additional reasoning or evidence to support your explanation.} \\
            }
        }
    }
     \\ \\

    \theutterance \stepcounter{utterance}  
    & \multicolumn{4}{p{0.6\linewidth}}{
        \cellcolor[rgb]{0.561,0.737,0.561}{
            \makecell[{{p{\linewidth}}}]{
                \texttt{\tiny{[P1$\rangle$GM]}}
                \texttt{No} \\
\\ 
\texttt{**Reasoning:**} \\
\\ 
\texttt{To determine the answer, we must analyze the causal chain and the specific counterfactual scenario provided in the hypothetical world.} \\
\\ 
\texttt{1.  **Define the Causal Structure:**} \\
\texttt{    *   \$Zuph \textbackslash{}rightarrow Wibl\$} \\
\texttt{    *   \$Wibl \textbackslash{}rightarrow Uvzi\$} \\
\texttt{    *   The problem states: "zuph causes not wibl." This implies that if \$Zuph\$ is present (True), then \$Wibl\$ is absent (False). Let's denote presence as 1 and absence as 0. So, \$Zuph=1 \textbackslash{}implies Wibl=0\$.} \\
\texttt{    *   The problem states: "wibl causes not uvzi." This implies that if \$Wibl\$ is present (True), then \$Uvzi\$ is absent (False). So, \$Wibl=1 \textbackslash{}implies Uvzi=0\$.} \\
\\ 
\texttt{2.  **Determine the State of the World under the Counterfactual:**} \\
\texttt{    *   The question asks: "Would an individual be not uvzi if zuph instead of not zuph?"} \\
\texttt{    *   This is a counterfactual query comparing the state where \$Zuph=1\$ (the intervention) against the baseline where \$Zuph=0\$ (not zuph). However, the phrasing "if zuph instead of not zuph" specifically} \\
            }
        }
    }
    & & \\ \\

    \theutterance \stepcounter{utterance}  
    & & & \multicolumn{2}{p{0.3\linewidth}}{
        \cellcolor[rgb]{0.9,0.9,0.9}{
            \makecell[{{p{\linewidth}}}]{
                \texttt{\tiny{[GM$|$GM]}}
                \texttt{No} \\
            }
        }
    }
    & & \\ \\

    \theutterance \stepcounter{utterance}  
    & & & \multicolumn{2}{p{0.3\linewidth}}{
        \cellcolor[rgb]{0.9,0.9,0.9}{
            \makecell[{{p{\linewidth}}}]{
                \texttt{\tiny{[GM$|$GM]}}
                \texttt{no} \\
            }
        }
    }
    & & \\ \\

    \theutterance \stepcounter{utterance}  
    & & & \multicolumn{2}{p{0.3\linewidth}}{
        \cellcolor[rgb]{0.9,0.9,0.9}{
            \makecell[{{p{\linewidth}}}]{
                \texttt{\tiny{[GM$|$GM]}}
                \texttt{game\_result = WIN} \\
            }
        }
    }
    & & \\ \\

\end{supertabular}
}

\end{document}
